\pagestyle{empty}
%
% Deckblatt
%
\begin{figure}[p]
\begin{tikzpicture}
	% Rahmen
	\draw[rounded corners=4mm] (151mm,213mm) -- (151mm,196mm) -- (0mm,196mm) -- (0mm,145mm) -- (151mm,145mm) -- (151mm,63mm);
	% F
	\draw (62mm,3mm) -- (62mm,59mm) -- (102mm,59mm) (62mm,31mm) -- (83mm,31mm);
	% A
	\draw (73mm,3mm) -- (106mm,58mm) -- (106mm,3mm) (90mm,31mm) -- (106mm,31mm);
	% U
	\draw (113mm,58mm) -- (113mm,18mm) arc (180:360:19mm)  -- (151mm,58mm);

	% FAU Logo
	\node at (25mm,175mm){\includegraphics[keepaspectratio=true,width=35mm]{deckblatt/FAU-Siegel.png}};
	% Arbeitsart
	\node at (25mm,151mm){\Large\thethesis};
	% Titel
	\node at (97mm,175mm){
		\parbox{10cm}{
			\centering \bfseries \itshape \Large \thesisTitleFormatted
		}
	};
	% Autor
	\node at (97mm,151mm){
			\bfseries \itshape \Large \myname
	};
	
	\node at (60mm,59mm)[below left,inner sep=0]{
		\parbox{6cm}{
			\begin{flushright}
				\large Lehrstuhl für Informatik 6 \\ (Datenmanagement)
			\end{flushright}
		}
	};
	\node at (60mm,25mm)[above left]{
		\parbox{6cm}{
			\begin{flushright}
				\large Department Informatik
			\end{flushright}
		}
	};
	\node at (60mm,-2mm)[above left,inner sep=0]{
		\parbox{6cm}{
			\begin{flushright}
				\large Friedrich Alexander-\\Universität\\Erlangen-Nürnberg
			\end{flushright}
		}
	};
\end{tikzpicture}
\end{figure}
%
% Titelseite
%
\begin{titlepage}

  \begin{center}
    
    {\Huge \bfseries
      \thesistitle{}\\
    } 
    
    \vspace*{1cm}
    \thethesis{} im Fach Informatik
    \vspace{2cm}
    
    {\large vorgelegt von} \\
    \vspace*{0.7cm}
    {\Large \bfseries \myname} \\
    \vspace*{0.7cm}
    {\large geb. \birthday{} in \birthplace{}} 
    
    \vspace{1.5cm}
    
    angefertigt am 

    \vspace{1cm}
    
    {\bfseries
      Lehrstuhl für Informatik 6 (Datenmanagement) \\
      Department Informatik \\
      Friedrich-Alexander-Universität Erlangen-Nürnberg (FAU)\\
      }
    
    \vspace{0.5cm}
\end{center}
\begin{tabbing}
    Betreuer: \= \corrector{}\\
    \> \tutor{} 
\end{tabbing}
    \vspace{0.25cm}
    
\begin{tabbing}
  Beginn der Arbeit: \startofwork{} \\
  Abgabe der Arbeit: \dayofdoom{}
\end{tabbing}

\end{titlepage}

\vspace{0,5cm}
%
% Leere Seite für doppelseitiges empty Layout einfügen
\cleardoubleemptypage
%
% Seitenlayout mit Seitennummern, allerdings kein Header oder Footer
\pagestyle{plain}
%
% Römische Seitenzahlen
\pagenumbering{roman}
% Da wir 6 Seiten ohne Seitennummerierung im Seitenlayout empty verwenden, reicht ein 
% \usepackage[pdftex,pdfpagelabels,plainpages=false,]{hyperref} nicht mehr aus. Es muss zusätzlich noch der 
% Seitenzähler auf 7 anstatt wie defaultmäßig auf 1 für die folgend Seiten gesetzt werden.
\setcounter{page}{7}
%
\chapter*{Kurzfassung}
\section*{\thesistitle}
\selectlanguage{ngerman}

Das Open Source Lernmanagement System ILIAS ist ein umfangreiches Werkzeug, um Studenten an Universitäten beim Lernen zu unterstützen. Eine der großen Funktionen von ILIAS ist das Test Modul. Mit Hilfe dieses Moduls können Dozenten Tests erstellen, die ihre Studenten beispielsweise lösen müssen, um wichtige Dokumente freizuschalten. Natürlich bietet das Modul auch die Möglichkeit Prüfungen zu erstellen, die dann online stattfinden.

Während das Test Modul von ILIAS nativ bereits eine große Anzahl an verschiedenen Fragetypen unterstützt - wie zum Beispiel Single Choice, Multiple Choice und Sortierungsfragen -, müssen komplexere Fragetypen nachträglich mit Hilfe von Plugins hinzugefügt werden.

In diesem Projekt - genannt \eigenname{\textbf{Q}uestion \textbf{P}lugin for \textbf{I}lias - \textbf{SQL} \textbf{(QPI-SQL)}} - wurde ein ILIAS Plugin - genannt \eigenname{assSQLQuestion} - für diesen Zweck entwickelt. Das Plugin ermöglicht Dozenten interaktive Fragen zum Themenbereich SQL zu stellen. Studenten können dabei ihre Antworten vor der Abgabe ausführen. Die Abgaben werden dabei außerdem automatisch - durch das Anwenden verschiedener Metriken - bewertet. 
%
% Leere Seite für doppelseitiges plain Layout einfügen
\cleardoubleplainpage
%
% Inhaltsverzeichnis
\tableofcontents
%
% Leere Seite für doppelseitiges plain Layout einfügen
\cleardoubleplainpage
%
% Seitenlayout mit Seitennummern, Header und Footer
\pagestyle{headings}
%
% Arabische Seitenzahlen
\pagenumbering{arabic}
%
% Seitenzähler auf 1 setzen, das benötigt das Paket hyperref
\setcounter{page}{1}
